
\chapter{MARCO TEÓRICO}
\label{chap:marco_teorico}

%% ------------------------------------------------------------------
\section{Biodeterioro microbiano en sustratos graníticos}
\label{sec:biodeterioro}

El granito es una roca ígnea intrusiva de textura fanerítica compuesta
principalmente por cuarzo (SiO$_2$), feldespatos potásicos (KAlSi$_3$O$_8$),
plagioclasas (NaAlSi$_3$O$_8$--CaAl$_2$Si$_2$O$_8$) y micas (biotita,
moscovita)~\cite{mnaahp2015}. Esta composición lo hace susceptible a la
biocorrosión química mediada por el microbioma colonizador, que actúa
a través de dos mecanismos principales: la acidólisis y la
quelatación iónica ~\cite{dakal2012, yadav2025}. En el primer
mecanismo, bacterias heterótrofas y hongos secretan ácidos orgánicos de
bajo peso molecular, incluyendo ácido oxálico, cítrico, glucónico,
láctico y 2-cetoglucónico, que reducen el pH local del microambiente
de la biopelícula y disuelven activamente los feldespatos y silicatos
de la matriz mineral~\cite{yadav2025, gutarowska2015}. En el segundo, estos
mismos ácidos actúan como agentes quelantes que secuestran cationes
estructurales esenciales (Fe$^{2+}$, Al$^{3+}$, Mg$^{2+}$, Ca$^{2+}$) de
la red cristalina, debilitando la cohesión mineral~\cite{yadav2025}.
Géneros bacterianos como \textit{Bacillus}, \textit{Arthrobacter} y
\textit{Streptomyces} han sido identificados como productores de ácido
2-cetoglucónico, un quelante altamente eficiente de la sílice presente
en monumentos graníticos~\cite{gutarowska2015}.

Adicionalmente, bacterias quimiolitoautotróficas del ciclo del nitrógeno
(\textit{Nitrosomonas}, \textit{Nitrospira}) y del azufre
(\textit{Thiobacillus}, \textit{Sulfurimonas}) producen ácido nítrico y
sulfúrico como subproductos metabólicos, que corroen directamente los
minerales carbonatados y silicatados de la roca~\cite{dakal2012,
suchy2024,ding2020}. A nivel físico, las hifas fúngicas ejercen presión mecánica
de turgencia que les permite penetrar microporos y
fisuras en proceso casmolítico y criptoendolítico, ampliando
progresivamente el sistema de poros del granito y facilitando la
entrada de agua~\cite{yadav2025, gorbushina2007}.

%% ------------------------------------------------------------------
\section{Biopelículas subaéreas en patrimonio pétreo}
\label{sec:biopelículas}

Las biopelículas subaéreas (SABs, \textit{Subaerial Biofilms})
constituyen la forma dominante de vida microbiana en superficies
líticas expuestas a la atmósfera~\cite{villa2016}. Su formación sigue
un proceso secuencial regulado molecularmente: (1) adhesión reversible
mediada por fuerzas de Van der Waals y electrostáticas; (2) adhesión
irreversible mediante proteínas de adhesión de superficie, pili tipo IV
y fimbrias; (3) maduración con producción de sustancias poliméricas
extracelulares (EPS); y (4) dispersión de células al
entorno~\cite{flemming2010}. La matriz de EPS, que puede representar más del 90\,\% del volumen de la biopelícula madura, es producida
por operones específicos como \textit{epsA-O} en
\textit{Bacillus subtilis}, los clusters \textit{pel} y \textit{psl} en
\textit{Pseudomonas aeruginosa}, y el operón \textit{csgA} responsable
de la síntesis de fibras de curli en
enterobacterias~\cite{flemming2010}. Esta matriz actúa como una esponja
biológica que retiene agua, captura nutrientes orgánicos del aire y
protege a las células de la desecación y los
biocidas~\cite{flemming2010, cappitelli2020}.

En el contexto del patrimonio pétreo, Villa et al.~\cite{villa2016}
proponen un marco ecológico fundamental para este proyecto: las
biopelículas subaéreas tienen un rol dual. Por un lado, pueden
inducir biodeterioro activo mediante la producción de ácidos, pigmentos
y presión mecánica; por el otro, pueden actuar como capa de
bioprotección que amortigua las fluctuaciones térmicas e hídricas,
consolida la superficie y limita la disolución química al entrar en
contacto con agua de lluvia o condensación~\cite{villa2016,
cappitelli2020}. Esta dualidad funcional implica que la sola
identificación taxonómica de microorganismos es insuficiente para
determinar el riesgo real de deterioro; es necesaria la caracterización
del repertorio funcional, específicamente de los genes accesorios,
para discernir qué mecanismos se encuentran genómicamente codificados
en cada linaje~\cite{cappitelli2020, skipper2022}.

Las estrategias de supervivencia en el nicho oligotrófico del granito
incluyen la esporulación (operón \textit{spo0A}/\textit{spoIIA}
/\textit{spoIIE} en Firmicutes, altamente prevalentes en ambientes
líticos), la melanización en hongos meristemáticos negros
(\textit{Ochroconis}, \textit{Coniosporium}) como escudo contra estrés
oxidativo y radiación UV, y la síntesis de carotenoides en
arqueas y bacterias halotolerantes como \textit{Rubrobacter}
spp.~\cite{zanardini2019, skipper2022, alonso2025, sterflinger2013,Galperin2022,branysova2022}. Estos rasgos
de supervivencia constituyen los biomarcadores genómicos de mayor
relevancia para el modelo predictivo de riesgo.

%% ------------------------------------------------------------------
\section{Metagenómica \textit{shotgun} y ensamblaje de genomas
  metagenómicos (MAGs)}
\label{sec:metagenomics}

La metagenómica \textit{shotgun} permite el acceso directo al material
genómico total de una comunidad microbiana sin necesidad de cultivo
previo, superando el sesgo de cultivabilidad que limita los estudios
clásicos al 1--5\,\% de la diversidad real~\cite{skipper2022,li2018deterioration, PitaGaleana2025}. En este enfoque, el ADN extraído de la muestra
ambiental se fragmenta y secuencia masivamente, generando millones de
lecturas cortas (\textit{reads}) que representan fragmentos aleatorios
de todos los genomas presentes en la comunidad~\cite{branysova2022,PitaGaleana2025}.

El proceso de reconstrucción de MAGs sigue tres etapas principales:
ensamblaje, binning y refinamiento. En la
etapa de ensamblaje, los \textit{reads} se ensamblan \textit{de novo} en
contigs mediante ensambladores basados en grafos de De~Bruijn, como
MEGAHIT~\cite{megahit2015} o SPAdes~\cite{metaspades2017}, que
fragmentan las lecturas en \textit{k-mers} y reconstruyen las secuencias
por solapamiento. En la etapa de \textit{binning}, los contigs se
agrupan en MAGs candidatos mediante dos señales complementarias: la
composición nucleotídica (frecuencia de tetra-nucleótidos, un
vector de 256 dimensiones) y el perfil de cobertura
diferencial (razón de \textit{reads} alineados sobre longitud del
contig a lo largo de múltiples muestras)~\cite{bowers2017,PitaGaleana2025}. La premisa biológica es que
contigs del mismo organismo tendrán firmas de composición y coberturas
similares. MetaBAT~2~\cite{metabat2019} modela esta relación con una
distribución de probabilidad Gaussiana, mientras que DAS~Tool~\cite{Sieber2018} integra
las predicciones de múltiples \textit{binners} para maximizar el
número de MAGs de alta calidad por muestra. El refinamiento manual
mediante Anvi'o~\cite{bowers2017, Eren2015} permite inspeccionar visualmente la coherencia de cada
\textit{bin} y eliminar contaminantes.

Esta aproximación supera directamente la limitación del 16S~rRNA
\textit{amplicon} identificada en la formulación del problema de la
tesis: mientras la secuenciación de amplicones provee información
taxonómica a nivel de género o especie, los MAGs reconstruyen el
genoma completo potencial de cada linaje, habilitando el
análisis pangenómico y la predicción del potencial
metabólico~\cite{PitaGaleana2025,dram2020,panaroo2020}.

%% ------------------------------------------------------------------
\section{Evaluación de calidad genómica}
\label{sec:quality}

La calidad de un MAG se define por dos métricas fundamentales:
completitud y contaminación~\cite{bowers2017}. La
completitud refleja qué fracción del genoma real ha sido capturada en
el MAG; la contaminación indica la proporción de contigs provenientes
de organismos distintos al organismo \textit{target}.

Ambas métricas se estiman contando los genes marcadores de
copia única (SCGs, \textit{Single-Copy Genes}): genes que evolucionaron
para existir exactamente en copia única en $>$97\,\% de todos los
procariotas. Un SCG ausente sugiere incompletitud; un SCG presente en
múltiples copias sugiere contaminación por contigs de otro organismo ~\cite{bowers2017,Parks2015}.

\begin{equation}
    \text{Completitud} =
    \frac{\text{SCGs detectados}}{\text{SCGs esperados para el linaje}}
    \times 100\%
    \label{eq:completeness}
\end{equation}

\begin{equation}
    \text{Contaminación} =
    \frac{\text{SCGs redundantes}}{\text{SCGs esperados}}
    \times 100\%
    \label{eq:contamination}
\end{equation}

CheckM2~\cite{checkm2} moderniza este cálculo mediante modelos de
\textit{machine learning} (\textit{gradient boosting} y redes
neuronales) entrenados sobre cientos de miles de genomas sintéticos simulados a partir de referencias de alta calidad. A diferencia de su predecesor, no requiere asignación
taxonómica previa y es más preciso en linajes divergentes como
Patescibacteria o el superphylum CPR~\cite{checkm2}.

El estándar internacional MIMAG~\cite{bowers2017} establece tres
categorías, resumidas en la tabla~\ref{tab:mimag}: (a) alta calidad
(HQ): completitud $\geq$90\,\%, contaminación $<$5\,\%, ARNr 16S
completo; (b) calidad media (MQ): 50--90\,\% completitud, $<$10\,\%
contaminación; (c) baja calidad: $<$50\,\%. El criterio del OE1
($>$70\,\%, $<$5\,\%) se posiciona como un umbral intermedio
conservador, apropiado para análisis pangenómicos que requieren
representación mínima del \textit{core genome}~\cite{bowers2017,
checkm2}.

\begin{table}[htbp]
  \centering
  \caption{Categorías de calidad MIMAG para MAGs bacterianos y
  arqueales~\cite{bowers2017}.}
  \label{tab:mimag}
  \resizebox{\textwidth}{!}{%
    \begin{tabular}{lccc}
      \toprule
      \textbf{Categoría} &
      \textbf{Completitud (\%)} &
      \textbf{Contaminación (\%)} &
      \textbf{Criterios adicionales} \\
      \midrule
      Alta calidad (HQ)  & $\geq 90$ & $< 5$      & ARNr + $\geq 18$ ARNt \\
      Media calidad (MQ) & $\geq 50$ & $< 10$     & ---                    \\
      Baja calidad (LQ)  & $< 50$    & Cualquiera & Solo referencial       \\
      \bottomrule
    \end{tabular}%
  }
\end{table}

%% ------------------------------------------------------------------
\section{Posicionamiento filogenómico y ANI}
\label{sec:phylogenomics}

La identidad nucleotídica promedio (ANI, \textit{Average Nucleotide
Identity}) es la métrica estándar actual para la delimitación de
especies procarióticas a nivel genómico, sustituyendo empíricamente a
la hibridación ADN-ADN (DDH) con mayor reproducibilidad y escala
computacional~\cite{fastANI2018}. El ANI se calcula como el promedio
ponderado de identidad de secuencia de todos los fragmentos ortólogos
entre dos genomas:

\begin{equation}
  \mathrm{ANI}(A,B) \;=\;
  \frac{\displaystyle\sum_{i}\, \mathrm{identity}(f_i^A, f_i^B) \times |f_i|}
       {\displaystyle\sum_{i}\, |f_i|}
  \label{eq:ani}
\end{equation}

\noindent donde $f_i^A$ y $f_i^B$ son fragmentos ortólogos de 3\,kb.
Un ANI $\geq$95\,\% (con fracción de alineamiento, AF $\geq$65\,\%)
equivale a un DDH $\geq$70\,\%, el umbral convencional de
especie~\cite{fastANI2018, gtdbtk2019}. FastANI~\cite{fastANI2018} implementa este
cálculo mediante \textit{locality-sensitive hashing} sobre
\textit{k-mers} (\textit{sketching} de MinHash), permitiendo
comparaciones a escala masiva.

El clasificador GTDB-Tk~\cite{gtdbtk2019} ubica los MAGs en un árbol
filogenómico de referencia usando 120 genes marcadores bacterianos
(\textit{bac120}) o 122 arqueales (\textit{arc122}), colocando cada
genoma en la taxonomía del \textit{Genome Taxonomy Database}
(GTDB)~\cite{gtdb2022}. GTDB corrige las inconsistencias del NCBI
normalizando los rangos taxonómicos por divergencia evolutiva relativa 
(RED, por sus siglas en inglés), siendo
especialmente relevante para MAGs ambientales que con frecuencia
representan linajes sin cultivos de referencia~\cite{gtdbtk2019,
gtdb2022}.

Un problema adicional al de la sola delimitación taxonómica surge
cuando se aplica el ANI a colecciones completas de MAGs recuperados
por metagenómica \textit{shotgun}: dado que el ensamblaje y el
\textit{binning} se realizan de forma independiente por muestra, un
mismo linaje presente en varios puntos de muestreo puede recuperarse
varias veces como genomas ligeramente distintos entre sí, debido a
diferencias en la profundidad de secuenciación o en el proceso de
ensamblaje de cada muestra. A este proceso de identificar y colapsar
genomas redundantes se le conoce como \textit{desreplicación}, y la
herramienta estándar para realizarlo es dRep \cite{olm2017}, que
agrupa los genomas de una colección en clústeres según su ANI y
selecciona, dentro de cada clúster, un único genoma representativo.

Internamente, dRep realiza esta comparación en dos pasos. Primero
ejecuta una comparación rápida entre todos los genomas mediante
\textit{MinHash} (Mash), que agrupa genomas claramente similares en
clústeres preliminares de forma poco costosa computacionalmente.
Luego, dentro de cada clúster preliminar, calcula el ANI de forma
precisa mediante ANIm o FastANI para confirmar la agrupación final.
El genoma representativo de cada clúster se escoge según un
\textit{score} de calidad que pondera completitud, contaminación y
contigüidad (N50), con criterios análogos a los de calidad descritos
en la sección 2.4.

El umbral de ANI usado en el segundo paso define qué tan estricta es
la agrupación: un umbral cercano al 95\,\% agrupa genomas de la
misma especie, de forma consistente con el umbral de DDH descrito
previamente, mientras que un umbral cercano al 99\,\% agrupa solo
genomas de la misma cepa. Este último valor representa además el
límite práctico de detección del método, ya que dos copias idénticas
de un mismo genoma comparadas entre sí suelen rendir valores de ANI
de aproximadamente 99.99\,\%, por lo que exigir una identidad mayor
ya no aporta resolución real \cite{olm2017, fastANI2018}. De este
modo, la desreplicación cumple dos funciones: elimina la redundancia
técnica generada por el ensamblaje independiente de cada muestra, y
define operativamente la unidad de análisis, el linaje, sobre la
cual se construye posteriormente el pangenoma comparativo descrito
en la siguiente sección.

IQ-TREE \cite{minh2020} es un programa de inferencia filogenética que reconstruye árboles mediante el criterio de máxima verosimilitud. Su algoritmo de búsqueda combina perturbaciones de tipo \textit{nearest neighbor interchange} estocástico con una estrategia de hill-climbing eficiente, lo que le permite explorar un espacio amplio de topologías sin caer prematuramente en óptimos locales, un problema frecuente en métodos de búsqueda heurística más simples. Antes de la búsqueda del árbol, IQ-TREE incorpora ModelFinder \cite{kalyaanamoorthy2017modelfinder}, un procedimiento que evalúa de forma automática el ajuste estadístico de distintos modelos de sustitución de aminoácidos o nucleótidos sobre el alineamiento de entrada y selecciona el más adecuado según criterios de información (AIC, BIC), evitando que el usuario deba fijar el modelo de forma arbitraria. El soporte estadístico de cada rama del árbol resultante se estima mediante el ultrafast bootstrap (UFBoot) \cite{hoang2018ufboot}, una aproximación computacionalmente eficiente al remuestreo bootstrap tradicional que aproxima de forma no sesgada los valores de soporte con un costo varios órdenes de magnitud menor, permitiendo ejecutar miles de réplicas en una fracción del tiempo que requeriría el bootstrap clásico de Felsenstein. Por estas características, IQ-TREE es ampliamente utilizado en filogenómica bacteriana para construir árboles focalizados a partir de alineamientos concatenados de genes marcadores, como los conjuntos bac120 y ar122 empleados por el GTDB \cite{gtdbtk2019}, en escenarios donde se requiere no solo inferir la topología del árbol, sino cuantificar la confianza estadística de cada agrupamiento, como la monofilia de un linaje frente a sus genomas de referencia.

\subsection{La jerarquía taxonómica}  

La clasificación taxonómica organiza a los organismos en diferentes niveles jerárquicos según sus características y relaciones evolutivas. El \textbf{filo} agrupa organismos que comparten características estructurales y evolutivas generales, mientras que la \textbf{clase} reúne grupos más específicos dentro de un mismo filo. El \textbf{orden} subdivide las clases de acuerdo con características adicionales compartidas entre los organismos, y la \textbf{familia} agrupa géneros que presentan una relación evolutiva cercana. El \textbf{género} reúne especies estrechamente relacionadas y constituye la primera parte del nombre científico de un organismo. Finalmente, la \textbf{especie} representa el nivel taxonómico más específico y permite identificar grupos de organismos con características genéticas y biológicas altamente similares. En conjunto, estos niveles siguen una jerarquía de mayor a menor amplitud: \textbf{filo $\rightarrow$ clase $\rightarrow$ orden $\rightarrow$ familia $\rightarrow$ género $\rightarrow$ especie}.

%% ------------------------------------------------------------------
\section{Análisis pangenómico comparativo}
\label{sec:pangenomics}

El pangenoma de una especie o linaje bacteriano representa el
conjunto total de genes presentes en todos los miembros del
grupo~\cite{tettelin2005}. Introducido formalmente por Tettelin et
al.~\cite{tettelin2005} con una división inicial entre genes núcleo y accesorios, 
la genómica computacional moderna lo estructura algorítmicamente en tres 
compartimentos principales~\cite{Page2015, panaroo2020}:
(a) el \textbf{\textit{core genome}} ($G_c$): genes presentes en
$\geq$95\,\% de los genomas; codifican funciones esenciales
conservadas (metabolismo central, replicación, traducción);
(b) el \textbf{\textit{shell genome}}: genes en el 15--95\,\% de los
genomas; representan adaptaciones frecuentes en el nicho; y
(c) el \textbf{\textit{cloud genome}} (o \textit{dispensable}): genes
presentes en $<$15\,\%; incluyen genes singulares (\textit{singletons})
posiblemente adquiridos por HGT para nichos específicos.

Los genes accesorios son el producto directo de la
transferencia horizontal de genes (HGT, \textit{Horizontal
Gene Transfer}): el mecanismo por el cual las bacterias adquieren
genes de organismos filogenéticamente distantes mediante conjugación,
transformación o transducción~\cite{soucy2015}. Este proceso permite la
rápida adquisición de rasgos fenotípicos, incluyendo la capacidad de
producir ácidos orgánicos o \textit{biopelículas}, sin necesidad de
mutaciones secuenciales, lo que lo convierte en el principal motor de
adaptación ecológica en procariotas~\cite{soucy2015}. En este microambiente, 
la propia matriz de la \textit{biopelícula} actúa como un reservorio de ADN 
extracelular que facilita y potencia estos eventos de transferencia 
génica~\cite{flemming2010}. Herramientas recientes como PPanGGOLiN~\cite{ppanggolin2020} han introducido modelos estadísticos basados en grafos para particionar 
el pangenoma en sus componentes de manera probabilística. Sin embargo, 
en colecciones de MAGs, la fragmentación intrínseca de los ensamblajes 
sigue provocando una sobreestimación de genes accesorios. Para superar 
esta limitación, Panaroo~\cite{panaroo2020} emplea una representación 
gráfica del pangenoma que corrige activamente los errores de anotación 
y ensamblaje mediante la propagación del contexto genómico vecino, 
lo que reduce de manera drástica la inflación artificial del genoma 
accesorio debida a la fragmentación.

%% ------------------------------------------------------------------
\section{Anotación funcional y predicción del potencial metabólico}
\label{sec:annotation}

La anotación funcional de MAGs sigue un proceso jerárquico 
descrito en la metagenómica computacional~\cite{PitaGaleana2025}, 
el cual inicia con la predicción de ORFs (\textit{Open Reading Frames}) 
mediante algoritmos como Prodigal~\cite{Hyatt2010} (frecuentemente 
integrado en \textit{pipelines} como Prokka~\cite{prokka2014}). 

Bakta~\cite{schwengers2021} constituye una evolución de este enfoque
de anotación estructural y funcional. A diferencia de \textit{pipelines}
anteriores como Prokka, que dependen de alineamientos de secuencias
contra bases de datos de referencia limitadas, Bakta incorpora un
módulo de identificación de secuencias sin alineamiento
(\textit{alignment-free sequence identification}, AFSI), el cual
calcula \textit{hashes} MD5 a partir de secuencias proteicas completas
para acelerar la asignación de homología. Al emplear una base de datos
taxonómicamente independiente, esta característica resulta
especialmente relevante para MAGs ambientales, donde la especie o el
género del organismo colonizador con frecuencia pertenece a linajes
raros o carece de representación cercana en las bases de datos
curadas. Para las secuencias no identificadas por AFSI, Bakta integra
fuentes complementarias como RefSeq y PSC (\textit{Protein Sequence
Clusters} derivados de UniRef90), así como modelos de covarianza de
Rfam para ARN no codificante. 

A continuación, se realiza una búsqueda de homología contra bases 
de datos de referencia y la posterior asignación a categorías 
funcionales, para culminar finalmente con la reconstrucción de 
rutas metabólicas~\cite{dram2020}.

Las principales bases de datos funcionales utilizadas en este tipo de
análisis son:

\begin{itemize}
  \item \textbf{KEGG Orthology (KO)}~\cite{kegg2023}: organiza genes
    en módulos metabólicos y rutas de señalización; la completitud de
    una ruta se calcula como la fracción de KOs esenciales detectados
    en el genoma.
  \item \textbf{COG} (\textit{Clusters of Orthologous Groups})~\cite{Galperin2022}:
    categorías funcionales de alto nivel (metabolismo, transporte,
    regulación de transcripción).
  \item \textbf{dbCAN/CAZy}~\cite{cantarel2009,zhang2018}: enzimas activas sobre carbohidratos,
    relevantes para la degradación de EPS y polisacáridos de la
    biopelícula.
  \item \textbf{CARD} (\textit{Comprehensive Antibiotic Resistance
    Database})~\cite{Alcock2020}: genes de resistencia antimicrobiana, útil para evaluar
    persistencia bajo tratamientos biocidas.
\end{itemize}

DRAM (\textit{Distilled and Refined Annotation of
Metabolism})~\cite{dram2020} integra estas bases de datos en un sistema
de anotación que produce tablas de completitud por módulo metabólico y
resúmenes de los principales ciclos biogeoquímicos (nitrógeno, azufre y 
carbono), así como del inventario de enzimas activas sobre 
carbohidratos (CAZymes)~\cite{dram2020}. Esta herramienta es especialmente
apropiada para el presente estudio dado que el mapeo de estas rutas 
metabólicas permite inferir el potencial de biodeterioro de la comunidad: 
los ciclos inorgánicos (N y S) están directamente vinculados a la disolución de 
minerales por acidólisis y biocorrosión~\cite{zanardini2019, dakal2012}, mientras 
que la síntesis de polisacáridos se relaciona con la producción de EPS para 
la biopelícula. De este modo, la destilación metabólica se alinea 
perfectamente con la identificación de factores de riesgo exigida en los 
objetivos del OE4.

La completitud de una vía metabólica se expresa formalmente en la
metodología de DRAM~\cite{dram2020} como:

\begin{equation}
    \text{Completitud}_{\text{vía}} =
    \frac{\text{módulos KO detectados}}
         {\text{módulos KO totales en la vía}}
    \times 100\%
    \label{eq:pathway}
\end{equation}

Los biomarcadores genómicos de mayor relevancia para el modelo
predictivo de riesgo ecológico incluyen: genes del ciclo TCA
(\textit{acnA}, \textit{sdhA}, \textit{fumA}) asociados a la producción de
ácidos orgánicos biocorrosivos~\cite{cappitelli2020, zanardini2019}; la
gluconato deshidrogenasa (\textit{gcd}) responsable de la síntesis
del ácido 2-cetoglucónico, un potente quelante de silicatos~\cite{gutarowska2015};
genes de biosíntesis de EPS (como el operón \textit{pel}) cruciales
para la arquitectura de la \textit{biopelícula}~\cite{flemming2010}; y los
operones de esporulación (con reguladores maestros como \textit{spo0A})
como marcadores de persistencia y resistencia extrema~\cite{Galperin2022}.

%% ------------------------------------------------------------------
\section{Limitaciones de la inferencia funcional basada en datos
  genómicos}
\label{sec:limitations}

La caracterización del potencial metabólico mediante análisis genómico
de MAGs representa una aproximación de alto poder inferencial, pero
con una limitación intrínseca en ausencia de datos de
metatranscriptómica. Dado que los MAGs analizados en esta tesis se
recuperaron a partir de la fracción cultivable (hisopo B), los
organismos representados demostraron ser metabólicamente viables al
proliferar bajo las condiciones de cultivo empleadas; sin embargo, el
propio cultivo constituye un filtro selectivo cuyas condiciones
(medio, temperatura, tiempo de incubación) no reproducen el ambiente
real del monumento (radiación UV, desecación, escasez de nutrientes).
Por ello, la detección de un gen no garantiza que su función se
exprese de la misma forma \textit{in situ}, y los resultados de este
proyecto describen el potencial metabólico codificado en los linajes
colonizadores, no la actividad biológica real sobre la piedra, tal
como se declara en el alcance de la tesis.
